
\section{Ishqalanish koeffitsiyenti}

\begin{enumerate}

\item YPX xodimi avtomobilning asfalt yo‘lidagi tormozlanish izi 40 m ga
teng ekanligini aniqladi. G‘ildiraklarning asfaltga ishqalanish
koeffitsiyenti 0,5 bo‘lsa, avtomobil qanday tezlik bilan (km/h da)
harakatlangan?

\begin{enumerate}[label=\Alph*)]
\item 54
\item 72
\item 90
\item 108
\end{enumerate}

\item 1 kg massali jism gorizontal tekislikda joylashgan. Jismga
2 N gorizontal kuch ta’sir qiladi. Ishqalanish koeffitsiyenti 0,3
bo‘lsa, ishqalanish kuchini aniqlang.

\begin{enumerate}[label=\Alph*)]
\item 1
\item 2
\item 3
\item 4
\end{enumerate}

\item 10 kg massali jism gorizontal tekislikda joylashgan.
Jismga gorizontga $30^\circ$ burchak ostida yo‘nalgan
50 N kuch ta’sir qiladi. Ishqalanish koeffitsiyenti 0,2 bo‘lsa,
ishqalanish kuchini aniqlang. $g=10\text{ m/s}^2$.

\begin{enumerate}[label=\Alph*)]
\item 10
\item 12
\item 15
\item 20
\end{enumerate}

\item Massasi 8 kg bo‘lgan jism o‘ziga qo‘yilgan kuch ta’sirida
stol bo‘ylab tekis harakatlanayotgan bo‘lsa, shu gorizontal
kuchning qiymati qanday? Ishqalanish koeffitsiyenti 0,3.
$g=10\text{ m/s}^2$.

\begin{enumerate}[label=\Alph*)]
\item 16
\item 20
\item 24
\item 30
\end{enumerate}

\item 3 kg massali brusok gorizontal prujina yordamida
gorizontal joylashgan doska bo‘ylab tekis tortilmoqda.
Agar prujina 5 cm ga uzaygan bo‘lsa, uning bikrligi qanday?
Brusok va doska orasidagi ishqalanish koeffitsiyenti 0,25.
$g=10\text{ m/s}^2$.

\begin{enumerate}[label=\Alph*)]
\item 100
\item 120
\item 150
\item 180
\end{enumerate}

\item 2 kg massali jism gorizontal kuch ta’sirida
gorizontal sirt bo‘ylab 2 m/s$^2$ tezlanish bilan
harakatlanmoqda. Agar jism va tekislik orasidagi
ishqalanish koeffitsiyenti 0,2 bo‘lsa,
shu kuchning miqdorini toping.
$g=10\text{ m/s}^2$.

\begin{enumerate}[label=\Alph*)]
\item 4
\item 6
\item 8
\item 10
\end{enumerate}

\item Yuk mashinasi o‘zining kuzovidagi mahkamlanmagan
yuk siljimasligi uchun qanday eng katta tezlanish bilan
harakatlanishi kerak? Yukning kuzov tubiga
ishqalanish koeffitsiyenti 0,2.
$g=10\text{ m/s}^2$.

\begin{enumerate}[label=\Alph*)]
\item 1
\item 2
\item 3
\item 4
\end{enumerate}

\item Odam 6 kg massali chanani arqon yordamida
doimiy tezlikda tortib bormoqda.
Bunda arqon gorizont bilan qandaydir burchak hosil qiladi,
va shu burchakning tangensi 0,75.
Chana va gorizontal sirt orasidagi
ishqalanish koeffitsiyenti 0,3.
Arqonning taranglik kuchini aniqlang.
$g=10\text{ m/s}^2$.

\begin{enumerate}[label=\Alph*)]
\item 12
\item 15
\item 18
\item 24
\end{enumerate}

\item 6 kg massali chanani ip yordamida gorizontga
$30^\circ$ burchak ostida yo‘nalgan 20 N kuch bilan
tortganda, u qanday tezlanish oladi?
Ishqalanish koeffitsiyenti 0,1.
$\sqrt3=1.7$.
$g=10\text{ m/s}^2$.

\begin{enumerate}[label=\Alph*)]
\item 1
\item 2
\item 3
\item 4
\end{enumerate}

\item 2,8 kg massali brusok vertical devor bo‘ylab
yuqoriga 70 N ga teng bo‘lgan va vertikalga burchak
ostida yo‘nalgan kuch yordamida ko‘chirilmoqda.
Agar $\sin\alpha=0.6$, hamda brusok va devor orasidagi
ishqalanish koeffitsiyenti 0,4 ekanligi ma’lum bo‘lsa,
brusokning tezlanishini toping.
$g=10\text{ m/s}^2$.

\begin{enumerate}[label=\Alph*)]
\item 2
\item 4
\item 6
\item 8
\end{enumerate}

\end{enumerate}



\begin{center}
\begin{tabular}{|c|c||c|c|}
\hline
1 & B & 6 & C \\
\hline
2 & B & 7 & B \\
\hline
3 & C & 8 & C \\
\hline
4 & C & 9 & B \\
\hline
5 & C & 10 & B \\
\hline
\end{tabular}
\end{center}











\section{Ishqalanishli qiya tekislik}

\begin{enumerate}

\item Jism qiya tekislik bo‘ylab tekis sirpanmoqda.
Agar jismning tekislikka ishqalanish koeffitsiyenti 0,2 bo‘lsa,
tekislikning gorizontga qiyalik burchagi kotangensi nimaga teng?

\begin{enumerate}[label=\Alph*)]
\item 5
\item 4
\item 2
\item 10
\end{enumerate}

\item Jism birinchi marta qiyalik burchagi $30^\circ$ bo‘lgan
qiya tekislik ustida, ikkinchi marta qiyalik burchagi
$60^\circ$ bo‘lgan qiya tekislik ustida qo‘yiladi.
Agar ikkala holda ham ishqalanish koeffitsiyenti 0,8 bo‘lsa,
birinchi holdagi ishqalanish kuchi ikkinchi holdagidan
necha foizga ko‘p?

\begin{enumerate}[label=\Alph*)]
\item 15
\item 20
\item 25
\item 30
\end{enumerate}

\item Jism balandligi 3 m va uzunligi 5 m bo‘lgan
qiya tekislikdan sirpanib tushadi.
Ishqalanish koeffitsiyenti 0,5.
Jismning tezlanishini toping.

\begin{enumerate}[label=\Alph*)]
\item 1
\item 2
\item 3
\item 4
\end{enumerate}

\item Jismga turtki bilan qiya tekislik bo‘ylab yuqoriga
yo‘nalgan tezlik berildi.
Agar qiya tekislikning balandligi 4 m,
uzunligi 5 m va ishqalanish koeffitsiyenti 0,5 bo‘lsa,
jism tezlanishining miqdorini toping.

\begin{enumerate}[label=\Alph*)]
\item 11
\item 9
\item 7
\item 13
\end{enumerate}

\item Jismga turtki bilan qiya tekislik bo‘ylab yuqoriga
yo‘nalgan tezlik berildi.
Qiya tekislikning balandligi 3 m, uzunligi 5 m,
va ishqalanish koeffitsiyenti 0,6.
Jismning yuqoriga harakatlanishidagi tezlanish
miqdori uning pastga harakatlanishidagi
tezlanish miqdoridan necha marta katta?

\begin{enumerate}[label=\Alph*)]
\item 3
\item 6
\item 12
\item 9
\end{enumerate}

\item Balandligi 5 m va qiyalik burchagi gorizontga
$45^\circ$ bo‘lgan qiya tekislik uchidan jism
sirpanishni boshlaydi.
Jismning qiya tekislikka ishqalanish koeffitsiyenti 0,19.
Jismning tushish oxiridagi tezligini toping.

\begin{enumerate}[label=\Alph*)]
\item 7
\item 8
\item 9
\item 10
\end{enumerate}

\item Jism balandligi 5 m va uzunligi 13 m
bo‘lgan qiya tekislikdan sirpanib tushadi.
Ishqalanish koeffitsiyenti 0,4.
Jismning qiya tekislik bo‘ylab
harakat vaqtini toping.

\begin{enumerate}[label=\Alph*)]
\item 10
\item 11
\item 12
\item 13
\end{enumerate}

\item Agar kichik shayba qiyalik burchagi
gorizontga $30^\circ$ bo‘lgan qiya tekislikdan
pastga tekis harakatlanib tushsa,
u xuddi shu moddadan tayyorlangan,
balandligi 2,5 m va qiyalik burchagi
$60^\circ$ bo‘lgan qiya tekislikdan
qancha vaqtda sirpanib tushadi?

\begin{enumerate}[label=\Alph*)]
\item 1
\item 2
\item 3
\item 4
\end{enumerate}

\item Jism qiya tekislikning uchigacha
chiqa olishi uchun unga qiya tekislik
bo‘ylab yuqoriga yo‘nalgan qanday minimal
boshlang‘ich tezlik berish kerak?
Qiya tekislikning balandligi 6 m,
uzunligi 10 m,
ishqalanish koeffitsiyenti 0,5.

\begin{enumerate}[label=\Alph*)]
\item 10
\item 12
\item 14
\item 16
\end{enumerate}

\item Uzunligi 26 m va balandligi 10 m
bo‘lgan qiya tekislik uchida
kichik brusok joylashgan.
Brusok va tekislik orasidagi
ishqalanish koeffitsiyenti 0,45.
Brusok tekislik asosigacha
yeta olishi uchun unga
qanday minimal tezlik berish kerak?

\begin{enumerate}[label=\Alph*)]
\item 2
\item 4
\item 6
\item 8
\end{enumerate}

\item Qiyalik burchagi gorizontga
$45^\circ$ bo‘lgan qiya tekislik bo‘ylab
shayba yuqoriga turtib yuborildi.
Qanchadir vaqtdan so‘ng u to‘xtaydi
va pastga sirpana boshlaydi.
Shaybaning tekislikka ishqalanish
koeffitsiyenti 0,8.
Shaybaning tushish vaqti
ko‘tarilish vaqtidan
necha marta katta?

\begin{enumerate}[label=\Alph*)]
\item 2
\item 3
\item 4
\item 5
\end{enumerate}

\item 5 kg massali jismni
balandligi 3 m va uzunligi 5 m
bo‘lgan qiya tekislik bo‘ylab
tekis ko‘tarishmoqda.
Bunda unga qiya tekislikka
parallel kuch qo‘yilgan.
Agar ishqalanish koeffitsiyenti
0,3 ga teng bo‘lsa,
shu kuchning miqdorini toping.

\begin{enumerate}[label=\Alph*)]
\item 32
\item 36
\item 42
\item 48
\end{enumerate}

\item Balandligi 3 m va uzunligi 5 m
bo‘lgan qiya tekislikda
50 kg massali yuk turibdi.
Yukni $1\,\text{m/s}^2$ tezlanish
bilan sudrab chiqarish uchun
unga tekislik bo‘ylab yo‘nalgan
qanday kuch qo‘yish kerak?
Ishqalanish koeffitsiyenti 0,2.

\begin{enumerate}[label=\Alph*)]
\item 430
\item 380
\item 330
\item 480
\end{enumerate}

\item Jismga gorizontal kuch qo‘ygan holda
qiya tekislik bo‘ylab yuqoriga ko‘tarishmoqda.
Bu kuchning miqdori jismga ta’sir qiluvchi
og‘irlik kuchidan ikki marta katta.
Qiya tekislikning balandligi 3 m,
uzunligi 5 m.
Agar ishqalanish koeffitsiyenti
0,2 ga teng bo‘lsa,
jismning tezlanishini toping.

\begin{enumerate}[label=\Alph*)]
\item 2
\item 4
\item 6
\item 8
\end{enumerate}

\item Agar qiya tekislik ustida turgan jismni
o‘ziga bog‘langan ip yordamida
gorizontal yo‘nalishda
jismga ta’sir qiluvchi og‘irlik kuchidan
ikki marta kichik kuch bilan tortilsa,
u qiya tekislikdan qanday tezlanish
bilan tusha boshlaydi?
Qiya tekislikning balandligi 3 m,
uzunligi 5 m.
Ishqalanish koeffitsiyenti 0,8.

\begin{enumerate}[label=\Alph*)]
\item 4
\item 6
\item 8
\item 10
\end{enumerate}

\end{enumerate}



\begin{center}
\begin{tabular}{|c|c||c|c|}
\hline
1 & A & 9 & C \\
\hline
2 & C & 10 & B \\
\hline
3 & B & 11 & B \\
\hline
4 & A & 12 & C \\
\hline
5 & D & 13 & A \\
\hline
6 & C & 14 & C \\
\hline
7 & D & 15 & B \\
\hline
8 & A &  &  \\
\hline
\end{tabular}
\end{center}

